\section{Branching on missed conflicts to select resolving candidates}
\label{sec:conflict-candidates}

The subset stream of Section~\ref{sec:euler-transversal} can inspect many
tetrahedra that never occur in a positive conflict.  A missed conflict itself
gives a smaller set of choices: every uniform resolving set must select at
least one of its tetrahedra.  The optional
\path{candidate_policy='conflicts'} uses these choices directly.  The
supplied-source default remains \path{candidate_policy='subsets'}; the
diagram driver exposes the new choice through
\path{transversal_candidate_policy='conflicts'} when
\path{sector_method='transversal'} is selected.

This change improves a parameterized search bound.  It retains the same
source model, polynomial LP policy, positive-ray recovery, and independently
checked negative tree format.  It establishes neither a small resolving
set for every knot exterior nor general quasi-polynomial recognition.

\subsection{The candidate oracle and the parameters}

Fix a source, or a supplied root Q mask, and use its positive extreme-ray
conflict hypergraph $\mathcal H$ from
Section~\ref{sec:euler-transversal}.  Write $h=\tau(\mathcal H)$ and let
$U$ be the tetrahedra with at least two allowed root Q types, with
$u=|U|\le t$.  Every edge of $\mathcal H$ is a nonempty subset of $U$.
Let $\delta$ be the maximum of $1$ and the sizes of all full conflict
sets encountered during this execution.  For an input-only bound it may
be replaced by
\[
 \delta(T)=\max\bigl(\{1\}\cup\{|D|:D\in\mathcal H\}\bigr).
\]
Thus $\delta\le\delta(T)\le\max(1,u)$.  The observed maximum is an
execution statistic; observing small conflicts does not bound $h$.

For a proposed $B\subseteq U$, the existing candidate oracle starts from
the unchanged root mask.  It explores a complete Euler tree but splits
only at conflict tetrahedra belonging to $B$.  It either returns a
compatible positive vector, completes a negative tree, or encounters an
incompatible positive root extreme ray whose full conflict set $D$ misses
$B$.  The last outcome refutes this candidate.  A work limit produces
\path{INCONCLUSIVE}, which is propagated rather than treated as a
refutation.  The root-extreme-ray premise follows from coordinate-face
inheritance and the basic-ray contract of
Section~\ref{sec:exact-polynomial-lp}.

The complete negative tree remains the mathematical exclusion certificate.
The learned conflicts and the candidate-selection history are search
metadata.  In particular, pruning a candidate-selection state does not
exclude any normal surface by itself.

\subsection{Iterative depth and a partition by first choice}

For each allowance $k=0,1,\ldots,u$, the implementation starts a fresh
depth-first candidate search at $(B,X)=(\varnothing,\varnothing)$.
The chosen and excluded sets satisfy $B\cap X=\varnothing$.  Such a
state represents possible transversals $R$ satisfying
\[
 B\subseteq R,\qquad R\cap X=\varnothing,\qquad |R|\le k.
\]
The excluded set applies only to candidate selection.  It never removes
Q columns or modifies the candidate oracle's root mask.

At a visited state, inspect the retained conflict cache.  If an edge
$D$ misses $B$, use it.  The implementation chooses the least available
missed edge, ordering first by $|D\setminus X|$ and then
lexicographically.  If no retained edge misses $B$, call the candidate
oracle once.  A positive or complete negative outcome finishes the source
query.  On refutation, its newly inserted missed edge supplies the next
choice at this same state; the state is not charged or queried again.

Put $D\setminus X=\{a_1<\cdots<a_s\}$.  If this set is empty, or if
$|B|=k$, prune the candidate-selection state.  Otherwise visit the children
\[
 \left(B\cup\{a_i\},\ X\cup\{a_1,\ldots,a_{i-1}\}\right),
 \qquad 1\le i\le s.
\]
Each child adds one previously absent tetrahedron.  The candidate oracle
uses only its chosen set $B$, not $X$.

\begin{lemma}[Coverage and uniqueness of the candidate partition]
\label{lem:conflict-candidate-partition}
Every transversal represented at a parent is represented at exactly one
of these children.  Within one depth allowance, distinct visited states
have distinct chosen sets $B$.
\end{lemma}
\begin{proof}
A represented transversal must meet $D$.  Since it avoids $X$ and $D$
misses $B$, it has a first selected element $a_i$ of $D\setminus X$.
The $i$th child retains that element and excludes precisely the earlier
available choices.  Every other child either requires an absent earlier
element or excludes $a_i$.  This proves the partition assertion.

For uniqueness, two distinct paths first diverge at children $i<j$ of
some state.  Every descendant of child $i$ retains $a_i$ in its chosen
set; every descendant of child $j$ excludes $a_i$ permanently.  Their
chosen sets cannot coincide.  Along a single path the chosen size strictly
increases, so an ancestor and descendant cannot coincide either.  Future
adaptive choices of missed edges do not alter these invariants.
\end{proof}

\subsection{A bounded cache may forget conflicts}

The new policy retains at most $256$ conflict sets in insertion order,
evicting the oldest when a new edge is inserted into a full cache.
Duplicates currently in the cache are not inserted again.  This cache
is separate from the existing bounded source-row dual cache.  The subset
comparison policy retains its original accumulating conflict list.

Eviction does not invalidate a candidate partition already made.  Its
chosen and excluded sets still record that partition, and the edge used
to form it remains a genuine member of $\mathcal H$.  At a later state,
forgetting an edge can cause an additional oracle call; that call must
either conclude with an independently replayable result or provide a
new actual refutation.  In the maintained generator, a refuting edge is
inserted last and the candidate oracle returns immediately.  Since the
cache has positive capacity, that edge remains available when the
suspended state resumes.  Consequently each visited state makes at most
one candidate-oracle call, including when earlier edges are relearned.

\begin{proposition}[Completion by the uniform transversal number]
\label{prop:conflict-candidate-completion}
With all operational limits removed, the conflict candidate policy
finishes by allowance $k=h$, for both positive and negative sources.
On a negative source its returned resolving set has minimum cardinality
$h$.  A positive result need not identify a transversal.
\end{proposition}
\begin{proof}
Fix a size-$h$ transversal $R$ of $\mathcal H$.  In the pass $k=h$,
the root represents $R$.  At every state representing $R$, any retained
or newly learned missed edge meets $R$ and therefore has an available
element outside $X$.  Lemma~\ref{lem:conflict-candidate-partition}
selects a child still representing $R$.  Such a state cannot be pruned
for exhausted allowance: if $|B|=h$ and $B\subseteq R$, then $B=R$,
so no edge misses it.  This argument uses only the validity of the edge
currently selected and survives arbitrary cache eviction.

The candidate tree for a true transversal cannot be refuted by an
incompatible positive root ray.  It has depth at most $|B|$, since a
split permanently resolves a previously unresolved member of $B$.
It therefore returns a compatible positive witness or a complete
negative tree.  Following the represented transversal proves termination
by pass $h$.

Conversely, a completed negative candidate tree implies that $B$ meets
every positive incompatible root ray.  A ray whose conflict set avoided
$B$ would retain a singleton choice at each split, follow a path through
that complete tree, and contradict its negative leaf.  Thus a negative
result has $h\le|B|\le k\le h$, proving minimum cardinality.
A compatible positive witness may instead be found before pass $h$;
its existence does not prove that the chosen set hits all incompatible
positive root rays.
\end{proof}

With forgetting, an individual queried candidate need not be a minimum
hitting set of the currently retained cache.  The minimum-cardinality
statement concerns a completed negative result, not every intermediate
proposal or an early positive result.

\subsection{A conflict-size parameter in the bit bound}

Let $L$ be the encoding length of the exact fixed-source linear data and
$A$ its number of anchor tuples.  A candidate with $i$ chosen tetrahedra
has at most
\[
 T_i=\frac{3^{i+1}-1}{2}
\]
Euler mask nodes.  Each candidate-state split has at most $\delta$
children.  The partition also prevents repeated chosen sets in one pass.

\begin{proposition}[Conflict-directed candidate bound]
\label{prop:conflict-candidate-bound}
Select the polynomial exact LP policy and disable operational limits.
The number $Q$ of attempted Euler mask nodes, including all unsuccessful
candidates and repeated depth passes, satisfies
\[
 Q\le
 \sum_{i=0}^{h}(h-i+1)
       \min\!\left\{\binom ui,\delta^i\right\}
       \frac{3^{i+1}-1}{2}
 <\frac{27}{8}(3\delta)^h.
\]
With the fixed conflict cache, total bit work is
\[
 O\!\left(A(3\delta)^h\operatorname{poly}(t,L)\right).
\]
\end{proposition}
\begin{proof}
In one pass, depth $i$ contains at most $\delta^i$ states because every
edge used for branching has at most that many available choices.
Lemma~\ref{lem:conflict-candidate-partition} also bounds this count by
$\binom ui$.  Every state makes at most one candidate-oracle call, of
cost at most $T_i$ mask nodes.  Passes $i,\ldots,h$ can revisit this
depth, giving the displayed sum.

For $a=3\delta\ge3$, discarding the binomial improvement gives
\[
 Q\le\frac32\sum_{i=0}^{h}(h-i+1)a^i
   <\frac32\frac{a^{h+2}}{(a-1)^2}
   \le\frac{27}{8}a^h.
\]
This also covers singleton conflicts and the no-conflict execution by
our convention $\delta\ge1$.  Repeated visits at larger depth allowances
are included, rather than assumed free.

Each mask node makes at most one envelope query and $A$ anchor queries.
Section~\ref{sec:exact-polynomial-lp} supplies their polynomial bit
cost and positive extreme-ray recovery.  Source algebra, bounded dual
reuse, geometry checks and certificate construction also use polynomial
work per mask node and per anchor.  Inspecting at most $256$ conflict
sets, sorting their available tetrahedra, and forming chosen/excluded
children costs polynomial work per candidate state.  The number of
states has the same upper sum with $T_i$ replaced by $1$, so this
bookkeeping fits the stated bound as well.
\end{proof}

The fixed conflict cache matters for this sharp expression.  Scanning
or deduplicating an unbounded learned-edge list can cost an additional
factor equal to the number of previously learned conflicts.  That
factor is not automatically polynomial in $t$.  The maintained bounded
cache avoids this obligation while preserving completeness.

For comparison, dropping the $\delta^i$ improvement and using $h\le u$
gives the crude bound $O(t4^t)$ mask nodes, including repeated passes.
The subset policy has the $O(4^t)$ crude count from
Section~\ref{sec:euler-transversal}.  The new policy is therefore an
optional comparison, not a universal domination claim.  When $\delta$
is bounded, the new parameter dependence removes the factor from
enumerating arbitrary tetrahedron subsets.  For example, bounded
$\delta$ and $h=O(\log t)$ would give polynomial search at an
authenticated one-vertex stage, where $A\le4t$ and $L$ is polynomial
in $t$.  No such structural bounds are established here.

\subsection{Limits, diagnostics, and independent acceptance}

The shared \path{max_hitting_subsets} option has a policy-specific unit.
For \path{subsets}, it counts each distinct subset inspected by the old
stream.  For \path{conflicts}, it counts visited chosen/excluded states,
including states pruned by cached edges and roots revisited at larger
allowances.  A state is charged before inspection; resuming that same
state after its one refuted candidate query incurs no second charge.
Every new child visit still requires an allowance.  A limit of zero
permits no candidate-state visit or candidate oracle query.

All Euler nodes, LP pivots, anchors, ellipsoid updates and feasibility
calls remain shared across candidate failures and depth passes.  Their
counts include discarded partial candidate trees.  Source algebra and
the dual cache remain shared during a supplied-source query.  A crush
starts new source-local algebra and conflict storage, while the diagram
driver passes the remaining global work allowances to its successor.
Exhaustion produces \path{INCONCLUSIVE}; caller exceptions propagate
with their original identity.

The counters \path{transversal_candidates} and
\path{transversal_restarts} distinguish actual candidate queries from
refutations.  In the new policy, \path{learned_conflicts} counts cache
insertions, including relearning after eviction; it is not a count of
globally distinct edges.  \path{maximum_conflict_size} records full
observed edge sizes, \path{peak_cached_conflicts} records the source
cache peak, and \path{conflict_cache_evictions} counts removals.  Across
crushes, work and eviction counts add, while sizes and peaks take maxima.

Only a completed candidate Euler tree supplies a negative packet.
Independent replay reconstructs its unchanged root mask and every child
assignment without trusting candidate order, cache history or resolving
metadata.  Positive results retain original normal-coordinate validation,
support extraction and authenticated crushing continuation.  The next
section tests a possible geometric route to small resolving sets; the
candidate bound itself provides no general quasi-polynomial theorem.
